\section{Objectif de la simulation}

La simulation a pour objectif de valider les algorithmes de perception, d'estimation, de planification et de controle sans exposer le vehicule reel a des risques materiels. Elle permet aussi de comparer des estimations avec une verite terrain disponible dans Gazebo.

Les objectifs techniques principaux sont:

\begin{itemize}
  \item fournir un modele de vehicule RC avec LiDAR, IMU et geometrie de base;
  \item simuler des pistes et scenarios reproductibles;
  \item connecter Gazebo a ROS 2 via des ponts de messages;
  \item reutiliser la pile APEX en mode simulation;
  \item enregistrer des runs et produire des donnees de diagnostic;
  \item tester differents modes d'estimation, de mapping et de raffinement offline.
\end{itemize}

La simulation n'est pas une preuve de fonctionnement reel. Elle sert surtout a accelerer le developpement, identifier des regressions et experimenter des scenarios difficiles a reproduire physiquement.

